\documentclass[11pt,a4paper]{article}

\usepackage[utf8]{inputenc}
\usepackage[T1]{fontenc}
\usepackage{amsmath,amssymb,amsthm,amsfonts}
\usepackage{geometry}
\geometry{margin=2.5cm}
\usepackage{hyperref}
\usepackage{microtype}
\usepackage{booktabs}

\newtheorem{theorem}{Theorem}[section]
\newtheorem{lemma}[theorem]{Lemma}
\newtheorem{proposition}[theorem]{Proposition}
\newtheorem{corollary}[theorem]{Corollary}
\theoremstyle{definition}
\newtheorem{definition}[theorem]{Definition}
\newtheorem{conjecture}[theorem]{Conjecture}
\theoremstyle{remark}
\newtheorem{remark}[theorem]{Remark}
\newtheorem{openproblem}[theorem]{Open Problem}

\title{\textbf{A Fisher--Rao Regularity Criterion for the 3D Incompressible Navier--Stokes Equations on the Dyadic Energy Simplex}}
\author{\textbf{Mariano Panzano Caballé} \\
\small Independent Researcher, Barcelona, Spain \\
\small Reproducibility \& Formal Artifacts: \texttt{DOI: 10.5281/zenodo.24137568}}
\date{\today}

\begin{document}

\maketitle

\begin{abstract}
We study the global regularity problem for the three-dimensional incompressible Navier--Stokes equations on $\mathbb{R}^3$ in the unforced regime ($\mathbf{f}=0$). Using Littlewood--Paley dyadic decomposition, Bony paraproduct commutator estimates, and the logarithmic Besov refinement of the Beale--Kato--Majda criterion due to Kozono and Taniuchi, we introduce an informational-geometric quantity on the simplex of dyadic energy fractions: the cumulative Fisher--Rao path length $\mathcal{L}_{FR}(T^*)$. We formulate a \emph{conjectural} regularity criterion stating that finiteness of $\mathcal{L}_{FR}(T^*)$ precludes finite-time singularity formation, and we prove a partial result under an additional non-concentration hypothesis (H2). We are explicit about what is proved, what is conjectural, and what remains open: in particular, we do \emph{not} claim an unconditional solution to the 3D global regularity problem, and we identify the precise analytic gap separating our partial result from such a solution.
\end{abstract}

\noindent\textbf{MSC 2020:} 35Q30, 76D03, 42B25, 53B12. \\
\textbf{Keywords:} Navier--Stokes equations, Fisher--Rao metric, Littlewood--Paley decomposition, Bony paraproducts, Kozono--Taniuchi criterion, conditional regularity.

\vspace{3mm}
\hrule
\vspace{4mm}

\section{Introduction and the Supercriticality Barrier}

Consider the Cauchy problem for the three-dimensional incompressible Navier--Stokes equations on $\mathbb{R}^3 \times (0, T)$ in the absence of external body forces ($\mathbf{f} = 0$):
\begin{equation}
\label{eq:ns_cauchy}
\begin{cases}
\partial_t \mathbf{u} + (\mathbf{u} \cdot \nabla)\mathbf{u} + \nabla p = \nu \Delta \mathbf{u}, & (x, t) \in \mathbb{R}^3 \times (0, T), \\
\nabla \cdot \mathbf{u} = 0, & (x, t) \in \mathbb{R}^3 \times (0, T), \\
\mathbf{u}(x, 0) = \mathbf{u}_0(x), & x \in \mathbb{R}^3,
\end{cases}
\end{equation}
where $\mathbf{u}: \mathbb{R}^3 \times [0, T) \to \mathbb{R}^3$ is the velocity field, $p: \mathbb{R}^3 \times [0, T) \to \mathbb{R}$ is the kinematic pressure, $\nu > 0$ is the kinematic viscosity, and $\mathbf{u}_0 \in \mathcal{S}(\mathbb{R}^3)$ is smooth and divergence-free.

Since Leray \cite{Leray1934}, global weak solutions are known to exist and satisfy the energy inequality:
\begin{equation}
\label{eq:leray_energy}
\|\mathbf{u}(t)\|_{L^2}^2 + 2\nu \int_0^t \|\nabla \mathbf{u}(s)\|_{L^2}^2 \, ds \le \|\mathbf{u}_0\|_{L^2}^2.
\end{equation}
However, \eqref{eq:leray_energy} is supercritical with respect to the scaling invariance $\mathbf{u}_\lambda(x, t) = \lambda \mathbf{u}(\lambda x, \lambda^2 t)$, $p_\lambda(x, t) = \lambda^2 p(\lambda x, \lambda^2 t)$, under which $\|\mathbf{u}_\lambda(\cdot, 0)\|_{L^2} = \lambda^{-1/2} \|\mathbf{u}_0\|_{L^2}$. Critical scale-invariant spaces follow the continuous embedding chain:
\begin{equation}
\dot{H}^{1/2}(\mathbb{R}^3) \hookrightarrow L^3(\mathbb{R}^3) \hookrightarrow \dot{B}^{-1 + 3/p}_{p, \infty}(\mathbb{R}^3) \hookrightarrow BMO^{-1}(\mathbb{R}^3) \hookrightarrow \dot{B}^{-1}_{\infty, \infty}(\mathbb{R}^3).
\end{equation}
Global regularity holds for small data in $BMO^{-1}$ (Koch--Tataru \cite{KochTataru2001}), and blow-up is excluded if the critical $L^\infty_t L^3_x$ norm remains bounded (Escauriaza--Seregin--\v{S}ver\'ak \cite{ESS2003}). Furthermore, Tao \cite{Tao2016} demonstrated that any regularity proof relying solely on energy identities and dyadic estimates that remain valid for averaged Navier--Stokes systems cannot prevent finite-time blow-up.

\paragraph{Relation to previous work.}
The use of Fisher--Rao geometry in fluid dynamics has appeared in the context of statistical mechanics of turbulence, but to our knowledge it has not been used as a regularity criterion. Our approach differs from Cheskidov--Shvydkoy \cite{CheskidovShvydkoy2010} and Kukavica--Vicol \cite{KukavicaVicol2009} in that we do not seek a new critical norm, but rather a geometric quantity on the spectral simplex that is finite for laminar and moderately turbulent flows and whose divergence would signal a cascade toward blow-up. The connection to Constantin--Fefferman \cite{ConstantinFefferman1993} is motivational: their geometric depletion criterion suggests that coherent vorticity alignment is the mechanism that could enforce (H1), but we do not establish this link.

---

\section{Littlewood--Paley Decomposition and Bony Paraproduct Estimates}

Let $\{\varphi_j\}_{j \in \mathbb{Z}}$ be a homogeneous dyadic partition of unity supported on annuli $\mathcal{C}_j = \{\xi \in \mathbb{R}^3 : 2^{j-1} \le |\xi| \le 2^{j+1}\}$ with $\sum_{j \in \mathbb{Z}} \varphi_j(\xi) = 1$ for $\xi \neq 0$. Define $\dot{\Delta}_j f = \mathcal{F}^{-1}(\varphi_j \hat{f})$ and $\dot{S}_j f = \sum_{k \le j-1} \dot{\Delta}_k f$, with homogeneous Besov norm $\|f\|_{\dot{B}^s_{p,r}} = \left\| \{2^{js} \|\dot{\Delta}_j f\|_{L^p}\}_{j \in \mathbb{Z}} \right\|_{\ell^r}$.

Taking the curl of \eqref{eq:ns_cauchy} yields the vorticity equation $\partial_t \boldsymbol{\omega} + (\mathbf{u} \cdot \nabla)\boldsymbol{\omega} - \nu \Delta \boldsymbol{\omega} = (\boldsymbol{\omega} \cdot \nabla)\mathbf{u}$, where $\boldsymbol{\omega} = \nabla \times \mathbf{u}$. Localizing to shell $j$ via Bony's paraproduct decomposition gives:
\begin{equation}
\partial_t \dot{\Delta}_j \boldsymbol{\omega} - \nu \Delta \dot{\Delta}_j \boldsymbol{\omega} = -\dot{S}_{j-1}\mathbf{u} \cdot \nabla \dot{\Delta}_j \boldsymbol{\omega} + \mathcal{R}_j(\mathbf{u}, \boldsymbol{\omega}) + \dot{\Delta}_j\big((\boldsymbol{\omega} \cdot \nabla)\mathbf{u}\big).
\end{equation}

\begin{lemma}[Dyadic Commutator Control]
\label{lem:bony_bound}
Let $\mathbf{u}$ be a smooth divergence-free solution to \eqref{eq:ns_cauchy} on $[0, T)$. For every $j \in \mathbb{Z}$, there exist universal constants $c_0, C_0 > 0$ such that:
\begin{equation}
\frac{1}{2}\frac{d}{dt}\|\dot{\Delta}_j \boldsymbol{\omega}\|_{L^2}^2 + c_0 \nu 2^{2j} \|\dot{\Delta}_j \boldsymbol{\omega}\|_{L^2}^2 \le C_0 \|\boldsymbol{\omega}\|_{\dot{B}^0_{\infty, \infty}} \|\dot{\Delta}_j \boldsymbol{\omega}\|_{L^2}^2 + C_0 \sum_{|k-j| \le 2} \|\dot{\Delta}_k \boldsymbol{\omega}\|_{L^2} \|\dot{\Delta}_j \boldsymbol{\omega}\|_{L^2} \|\dot{S}_{j-1}\nabla \mathbf{u}\|_{L^\infty}.
\end{equation}
\end{lemma}

---

\section{Logarithmic Besov Regularity and Fisher--Rao Spectral Geometry}

By the Beale--Kato--Majda theorem \cite{BKM1984} and its logarithmic Besov refinement by Kozono and Taniuchi \cite{KozonoTaniuchi2000}, a smooth solution $\mathbf{u}$ extends past $T^* < \infty$ whenever $\int_0^{T^*} \|\boldsymbol{\omega}(t)\|_{\dot{B}^0_{\infty, \infty}} \, dt < \infty$, via the logarithmic estimate:
\begin{equation}
\label{eq:kt_ineq}
\|\nabla \mathbf{u}\|_{L^\infty} \le C_1 \left( 1 + \|\boldsymbol{\omega}\|_{\dot{B}^0_{\infty, \infty}} \ln^+\!\left(\frac{\|\mathbf{u}\|_{H^s}}{\|\boldsymbol{\omega}\|_{\dot{B}^0_{\infty, \infty}}}\right) \right), \qquad s > 5/2.
\end{equation}

\begin{definition}[Dyadic Energy Simplex and Fisher--Rao Path Length]
Assume $\|\mathbf{u}(t)\|_{L^2} > 0$ on $[0,T^*)$. Define
\[
p_j(t) := \frac{\|\dot{\Delta}_j \mathbf{u}(t)\|_{L^2}^2}{\sum_{k}\|\dot{\Delta}_k \mathbf{u}(t)\|_{L^2}^2},
\qquad
q_j(t) := \sqrt{p_j(t)},
\qquad
\sum_j q_j(t)^2 = 1.
\]
The instantaneous Fisher--Rao speed and cumulative path length are
\[
\mathcal{I}_{FR}(t) := 2\Big(\sum_j (\partial_t q_j(t))^2\Big)^{1/2},
\qquad
\mathcal{L}_{FR}(T^*) := \int_0^{T^*} \mathcal{I}_{FR}(t)\,dt.
\]
The map $p_j \mapsto q_j$ is an isometry from the Fisher--Rao simplex onto the unit sphere $S^\infty \subset \ell^2(\mathbb{Z})$, and removes the coordinate singularity at $p_j = 0$.
\end{definition}

---

\section{A Partial Regularity Result and a Conjecture}

\subsection{Statement of the partial result}

\begin{theorem}[Partial Fisher--Rao regularity criterion]
\label{thm:partial}
Let $\mathbf{u}_0 \in H^s(\mathbb{R}^3)$ ($s>5/2$) be divergence-free, and let $\mathbf{u} \in C([0,T^*); H^s(\mathbb{R}^3))$ be the local smooth solution to \eqref{eq:ns_cauchy} with $\mathbf{f}=0$. Assume:
\begin{enumerate}
  \item[\emph{(H1)}] $\mathcal{L}_{FR}(T^*) := 2\int_0^{T^*}\|\partial_t q(t)\|_{\ell^2}\,dt < \infty$;
  \item[\emph{(H1$'$)}] $\displaystyle\lim_{J\to\infty}\sup_{t\in[0,T^*)}\sum_{j\ge J} p_j(t) = 0$;
  \item[\emph{(H2)}] There exist $J_0\in\mathbb{Z}$ and $g\in L^1(0,T^*)$ such that
        $\sup_{j\ge J_0} 2^{5j/2}\sqrt{p_j(t)} \le g(t)$ for all $t\in[0,T^*)$.
\end{enumerate}
Then $\int_0^{T^*}\|\boldsymbol{\omega}(t)\|_{\dot{B}^0_{\infty,\infty}}\,dt<\infty$, and $\mathbf{u}$ extends smoothly past $T^*$.
\end{theorem}

\begin{proof}
Under (H1), Cauchy--Schwarz on $\ell^2(\mathbb{Z})$ gives
\[
\sum_{j \in \mathbb{Z}} |\partial_t p_j(t)| = 2 \sum_{j \in \mathbb{Z}} q_j(t) |\partial_t q_j(t)| \le 2 \left(\sum_{j\in\mathbb{Z}} q_j(t)^2\right)^{1/2} \left(\sum_{j\in\mathbb{Z}} (\partial_t q_j(t))^2\right)^{1/2} = \mathcal{I}_{FR}(t),
\]
so $\int_0^{T^*} \sum_j |\partial_t p_j(t)|\,dt \le \mathcal{L}_{FR}(T^*) < \infty$ and each $p_j(t)$ admits a limit as $t \nearrow T^*$. Hypothesis (H1$'$) guarantees uniform non-escape of energy mass to arbitrarily high shells. By Bernstein's inequality,
\[
\|\dot{\Delta}_j\boldsymbol{\omega}\|_{L^\infty} \le C_B 2^{5j/2}\|\dot{\Delta}_j\mathbf{u}\|_{L^2} = C_B \|\mathbf{u}(t)\|_{L^2} 2^{5j/2}\sqrt{p_j(t)}.
\]
Splitting at $J_0$: the low-frequency part $j<J_0$ is controlled by $\|\mathbf{u}_0\|_{L^2}$ and the finite sum, while the high-frequency part $j\ge J_0$ is controlled by (H2). Hence $\|\boldsymbol{\omega}(t)\|_{\dot{B}^0_{\infty,\infty}} \le C(\|\mathbf{u}_0\|_{L^2} + g(t))$, which is integrable on $(0,T^*)$. The Kozono--Taniuchi criterion \eqref{eq:kt_ineq} then yields smooth extendability past $T^*$.
\end{proof}

\subsection{What is missing: the analytic gap}

The proof above uses (H2) as a black box. We now state precisely why (H2) is not derived from the equations, and what would be needed to remove it.

\begin{remark}[The analytic gap]
\label{rem:gap}
Hypothesis (H1) controls the total variation of the spectral distribution in $L^1_t$. Hypothesis (H1$'$) excludes mass escaping to arbitrarily high shells as $t \nearrow T^*$. Hypothesis (H2) is a uniform-in-time, $\sup$-in-$j$ bound on $2^{5j/2}\sqrt{p_j(t)}$. The gap between (H1)--(H1$'$) and the desired conclusion is exactly the gap between $L^1_t$ control of the total variation and $\sup_j$ control of the weighted amplitude. Neither an interpolation argument nor a direct PDE estimate is currently available: the weight $2^{5j/2}$ is not controlled uniformly in $j$ by the second moment $2^{2j}$ available from Leray's inequality, and a direct PDE estimate would essentially solve the 3D regularity problem.
\end{remark}

\subsection{The conjecture}

\begin{conjecture}[Fisher--Rao regularity criterion]
\label{conj:main}
Let $\mathbf{u}_0 \in H^s(\mathbb{R}^3)$ ($s>5/2$) be divergence-free, and let $\mathbf{u}$ be the local smooth solution to \eqref{eq:ns_cauchy} with $\mathbf{f}=0$. If $\mathcal{L}_{FR}(T^*) < \infty$ and (H1$'$) holds, then $\int_0^{T^*}\|\boldsymbol{\omega}(t)\|_{\dot{B}^0_{\infty,\infty}}\,dt < \infty$, and $\mathbf{u}$ extends smoothly past $T^*$.
\end{conjecture}

\begin{remark}[Status]
Conjecture~\ref{conj:main} is proved in Theorem~\ref{thm:partial} under the additional non-concentration hypothesis (H2). We do not know whether (H1) and (H1$'$) alone suffice. We do not know whether (H1) holds for all smooth unforced solutions; proving this would imply global regularity and is therefore at least as hard as the Clay Millennium Problem.
\end{remark}

---

\section{Scope, Evidence, and What Is Not Claimed}

\subsection{Three levels of statement}

We separate three levels to avoid any ambiguity:
\begin{enumerate}
  \item \textbf{Proved (Theorem~\ref{thm:partial}):} Under (H1)--(H1$'$)--(H2), no finite-time blow-up.
  \item \textbf{Conjectured (Conjecture~\ref{conj:main}):} Under (H1)--(H1$'$) alone, no finite-time blow-up.
  \item \textbf{Open:} Whether (H1) holds for all smooth unforced data. This is at least as hard as the Clay problem.
\end{enumerate}

\subsection{Numerical evidence}

We do not report DNS data in this manuscript. The companion LES paper reports spectral diagnostics for the Taylor--Green vortex at $\mathrm{Re}=1600$; those diagnostics are consistent with the finiteness of $\mathcal{L}_{FR}$ over the simulated window, but they are not a proof of (H1).\footnote{When DNS data become available, the discrete path length is evaluated via the exact spherical Fisher--Rao geodesic arc length $\Delta\mathcal{L}_{FR}(t_n,t_{n+1}) = 2\arccos\big(\min(1,\sum_j \sqrt{p_j(t_n)p_j(t_{n+1})})\big)$, which is free of finite-difference division singularities.} Numerical evidence for a hypothesis is not a proof of it, and we do not claim otherwise.

\subsection{What is not claimed}

\begin{itemize}
  \item We do not claim an unconditional solution to the 3D global regularity problem.
  \item We do not claim that $\kappa=\ln 10$ is an invariant of \eqref{eq:ns_cauchy}. It is a modeling hypothesis, used only in the LES companion paper, not in Theorem~\ref{thm:partial} or Conjecture~\ref{conj:main}.
  \item We do not claim that Fisher--Rao path length is invariant under Tao's averaged flow. We expect it is \emph{not}, since averaging destroys inter-shell phase coherence; but we do not prove this. Whether this suffices to place Conjecture~\ref{conj:main} outside Tao's obstruction is open.
  \item We do not claim sharpness of (H1), (H1$'$), or (H2). We do not know whether they are necessary, sufficient alone, or improvable.
\end{itemize}

---

\section{Open Problems}

We collect the precise questions that remain open, in decreasing order of ambition.

\begin{openproblem}[Unconditional regularity]
Prove or disprove Conjecture~\ref{conj:main} without hypothesis (H2). Equivalently: prove that (H1) and (H1$'$) alone imply $\int_0^{T^*}\|\boldsymbol{\omega}\|_{\dot{B}^0_{\infty,\infty}}dt<\infty$.
\end{openproblem}

\begin{openproblem}[Validity of (H1)]
Prove that every smooth unforced solution of \eqref{eq:ns_cauchy} on $\mathbb{R}^3$ satisfies $\mathcal{L}_{FR}(T^*)<\infty$ for all $T^*>0$ for which the solution exists. This would imply global regularity.
\end{openproblem}

\begin{openproblem}[Geometric mechanism]
By Constantin--Fefferman \cite{ConstantinFefferman1993}, vortex stretching is depleted when the unit vorticity field $\boldsymbol{\xi}=\boldsymbol{\omega}/|\boldsymbol{\omega}|$ satisfies a H\"older-type coherence condition. Is there a direct link between the Biot--Savart geometry of $\boldsymbol{\xi}$ and the Fisher--Rao speed $\mathcal{I}_{FR}(t)$? A positive answer would provide a mechanism for (H1).
\end{openproblem}

\begin{openproblem}[Relation to Tao's obstruction]
Tao \cite{Tao2016} rules out regularity proofs that remain valid under averaged flows. Is $\mathcal{L}_{FR}$ non-invariant under averaging? If so, does that suffice to place Conjecture~\ref{conj:main} outside the class of arguments excluded by Tao?
\end{openproblem}

\section*{Acknowledgments and Formal Verification Artifacts}
Algebraic identities for the Bhattacharyya--Hellinger sphere regularization ($p_j = q_j^2$) and simplex variation bounds were machine-checked in Lean~4 (\texttt{RiemannFisherRaoHellinger.lean} (over $\mathbb{R}$), \texttt{BonyCommutatorDyadic.lean}) and archived in the companion repository (\texttt{DOI: 10.5281/zenodo.24137568}).

\begin{thebibliography}{99}

\bibitem{Leray1934}
J.~Leray.
\newblock Sur le mouvement d'un liquide visqueux emplissant l'espace.
\newblock {\em Acta Mathematica}, 63(1):193--248, 1934.

\bibitem{BKM1984}
J.~T. Beale, T.~Kato, and A.~Majda.
\newblock Remarks on the breakdown of smooth solutions for the 3-{D} {E}uler equations.
\newblock {\em Communications in Mathematical Physics}, 94(1):61--66, 1984.

\bibitem{ConstantinFefferman1993}
P.~Constantin and C.~Fefferman.
\newblock Direction of vorticity and the problem of global regularity for the {N}avier--{S}tokes equations.
\newblock {\em Indiana University Mathematics Journal}, 42(3):775--789, 1993.

\bibitem{KukavicaVicol2009}
I.~Kukavica and V.~Vicol.
\newblock On the radius of analyticity of solutions to the three-dimensional {N}avier--{S}tokes equations.
\newblock {\em Proceedings of the American Mathematical Society}, 137(2):669--677, 2009.

\bibitem{CheskidovShvydkoy2010}
A.~Cheskidov and R.~Shvydkoy.
\newblock The regularity of weak solutions of the 3{D} {N}avier--{S}tokes equations in $B^{-1}_{\infty,\infty}$.
\newblock {\em Archive for Rational Mechanics and Analysis}, 195(1):159--169, 2010.

\bibitem{KozonoTaniuchi2000}
H.~Kozono and Y.~Taniuchi.
\newblock Bilinear estimates in {BMO} and the {N}avier--{S}tokes equations.
\newblock {\em Mathematische Zeitschrift}, 235(1):173--194, 2000.

\bibitem{KochTataru2001}
H.~Koch and D.~Tataru.
\newblock Well-posedness for the {N}avier--{S}tokes equations.
\newblock {\em Advances in Mathematics}, 157(1):22--35, 2001.

\bibitem{ESS2003}
L.~Escauriaza, G.~Seregin, and V.~\v{S}ver\'ak.
\newblock $L_{3,\infty}$-solutions of {N}avier--{S}tokes equations and backward uniqueness.
\newblock {\em Russian Mathematical Surveys}, 58(2):211--250, 2003.

\bibitem{Tao2016}
T.~Tao.
\newblock Finite time blowup for an averaged three-dimensional {N}avier--{S}tokes equation.
\newblock {\em Journal of the American Mathematical Society}, 29(3):601--674, 2016.

\end{thebibliography}

\end{document}
